\documentclass[10pt]{article}
\providecommand{\FIGDIR}{fig}
\newcommand{\DIGESTDATE}{29 July 2026}
% ===================== PM Research Watch - shared preamble =====================
\usepackage[T1]{fontenc}
\usepackage[utf8]{inputenc}
\usepackage{amsmath,amssymb}
\usepackage{tgpagella}
\usepackage{tgheros}
\renewcommand{\sfdefault}{qhv}
\renewcommand{\ttdefault}{lmtt}
\usepackage{microtype}
\usepackage[english]{babel}

\usepackage[a4paper,top=17mm,bottom=17mm,left=17mm,right=17mm,
            headsep=6mm,footskip=10mm]{geometry}
\usepackage{graphicx}
\usepackage{xcolor}
\usepackage{colortbl}
\usepackage{tikz}
\usetikzlibrary{calc}
\usepackage[most]{tcolorbox}
\usepackage{titlesec}
\usepackage{fancyhdr}
\usepackage{booktabs}
\usepackage{array}
\usepackage{longtable}
\usepackage{enumitem}
\usepackage{multicol}
\usepackage{ragged2e}
\usepackage[hidelinks]{hyperref}

% ------------------------------------------------------------------ palette
\definecolor{Ink}    {HTML}{14262E}
\definecolor{Deep}   {HTML}{0F4C5C}
\definecolor{Teal}   {HTML}{1F7A8C}
\definecolor{Sky}    {HTML}{4FA3C4}
\definecolor{Sage}   {HTML}{6FA287}
\definecolor{Amber}  {HTML}{E3A008}
\definecolor{Clay}   {HTML}{B5643C}
\definecolor{Coral}  {HTML}{D1495B}
\definecolor{Violet} {HTML}{6A5B9E}
\definecolor{Slate}  {HTML}{4A6572}
\definecolor{Paper}  {HTML}{FBF9F5}
\definecolor{Mist}   {HTML}{EEEAE2}
\definecolor{Rule}   {HTML}{D9D3C9}

\pagecolor{Paper}
\color{Ink}

\hypersetup{colorlinks=true, linkcolor=Deep, urlcolor=Teal, citecolor=Teal,
            pdftitle={PM Research Watch}, pdfauthor={Automated daily digest}}

% ------------------------------------------------------------------ headings
\titleformat{\section}
  {\sffamily\bfseries\large\color{Deep}}{}{0pt}
  {}[{\vspace{1.2mm}\color{Amber}\titlerule[1.1pt]}]
\titlespacing*{\section}{0pt}{5.5mm}{2.8mm}

\titleformat{\subsection}
  {\sffamily\bfseries\normalsize\color{Teal}}{}{0pt}{}
\titlespacing*{\subsection}{0pt}{3.6mm}{1.4mm}

\setlist[itemize]{leftmargin=4.2mm, itemsep=1.1mm, topsep=1.3mm, parsep=0pt,
                  label=\textcolor{Teal}{\footnotesize$\blacktriangleright$}}

% ------------------------------------------------------------------ header/footer
\pagestyle{fancy}
\fancyhf{}
\renewcommand{\headrulewidth}{0pt}
\renewcommand{\footrulewidth}{0pt}
\fancyhead[L]{\sffamily\scriptsize\color{Slate}PM RESEARCH WATCH\;\textcolor{Rule}{\textbar}\;\DIGESTDATE}
\fancyhead[R]{\sffamily\scriptsize\color{Slate}Particulate matter monitoring \& health effects}
\fancyfoot[C]{\sffamily\scriptsize\color{Slate}\thepage}

% ------------------------------------------------------------------ components
% metric tile
\newcommand{\metric}[3]{%
  \begin{tcolorbox}[enhanced, colback=white, colframe=Rule, boxrule=0.5pt,
    arc=1.6mm, left=2.2mm, right=2.2mm, top=1.8mm, bottom=1.8mm,
    borderline west={1.7mm}{0pt}{#3}, nobeforeafter, valign=center,
    height=17mm, width=\linewidth]
    \RaggedRight
    {\sffamily\bfseries\fontsize{17}{17}\selectfont\color{#3}#1}\\[0.5mm]
    {\sffamily\fontsize{6.5}{7.4}\selectfont\color{Slate}#2}
  \end{tcolorbox}}

% pull-quote / key insight box
\newtcolorbox{keybox}[1][]{enhanced, colback=Mist, colframe=Mist,
  arc=1.6mm, left=3.2mm, right=3.2mm, top=2.4mm, bottom=2.4mm,
  borderline west={1.7mm}{0pt}{Amber}, #1}

% subtopic band
\newcommand{\band}[2]{%
  \vspace{3.2mm}\noindent
  \begin{tcolorbox}[enhanced, colback=#2, colframe=#2, arc=1.2mm,
    left=2.6mm, right=2.6mm, top=1.1mm, bottom=1.1mm, boxrule=0pt,
    width=\linewidth, nobeforeafter]
    {\sffamily\bfseries\small\color{white}#1}
  \end{tcolorbox}\vspace{1.6mm}}

% one paper entry: {short cite}{title}{journal}{doi}{body}
\newcommand{\paperentry}[5]{%
  \par\vspace{1.9mm}\noindent
  {\sffamily\bfseries\footnotesize\color{Ink}#1}\;%
  {\sffamily\fontsize{7}{8}\selectfont\color{Slate}#3}\\[0.4mm]
  {\footnotesize\itshape\color{Deep}#2}\\[0.6mm]
  {\footnotesize #4}\\[0.3mm]
  {\fontsize{6.6}{7.6}\selectfont\color{Slate}\texttt{doi:}\,\href{https://doi.org/#5}{\texttt{\detokenize{#5}}}}
  \par}

% caption style
\newcommand{\figcap}[1]{\par\vspace{1.0mm}
  {\fontsize{7.2}{8.6}\selectfont\color{Slate}\textsf{\textbf{Fig.}}\; #1\par}\vspace{1.5mm}}

\providecommand{\permil}{\textperthousand}

\setlength{\parindent}{0pt}
\setlength{\parskip}{0pt}
\renewcommand{\arraystretch}{1.18}


\begin{document}

% ==================================================================== MASTHEAD
\thispagestyle{empty}
\begin{tikzpicture}[remember picture, overlay]
  \fill[Deep] (current page.north west) rectangle
      ($(current page.north east)+(0,-67mm)$);
  \fill[Amber] ($(current page.north west)+(0,-67mm)$) rectangle
      ($(current page.north east)+(0,-68.4mm)$);
  \foreach \x/\y/\r in {22/12/0.9, 41/26/1.5, 63/9/0.7, 88/31/1.1, 112/17/1.9,
                        138/29/0.8, 159/13/1.3, 178/54/1.0, 196/21/1.6,
                        30/44/0.6, 74/57/1.2, 125/48/0.9, 168/6/0.8, 191/59/0.7,
                        52/18/1.0, 100/60/0.6, 148/20/1.1, 66/49/1.4, 205/45/0.8}
    \fill[Sky, opacity=0.30] ($(current page.north west)+(\x mm,-\y mm)$) circle (\r mm);
\end{tikzpicture}

\vspace*{4mm}
{\sffamily\fontsize{7.6}{9}\selectfont\color{Sky}%
 AUTOMATED DAILY LITERATURE DIGEST\;\textbullet\;PhD RESEARCH WATCH\par}
\vspace{2.2mm}
{\sffamily\bfseries\fontsize{27}{29}\selectfont\color{white}%
 Particulate Matter\par}
{\sffamily\fontsize{27}{29}\selectfont\color{Sky}%
 Monitoring \& Health Effects\par}
\vspace{5.0mm}
{\sffamily\fontsize{8.4}{10}\selectfont\color{white}%
 Issue 2026-07-29\;\;\textcolor{Amber}{\textbar}\;\;PubMed entry window 29 July 2026\;\;%
 \textcolor{Amber}{\textbar}\;\;25 records in scope, 3 excluded\par}
\vspace{18mm}

% -------------------------------------------------------------- metric strip
\noindent
\begin{minipage}[t]{0.190\linewidth}\metric{25}{RECORDS IN SCOPE\\(3 excluded)}{Deep}\end{minipage}\hfill
\begin{minipage}[t]{0.190\linewidth}\metric{9}{SUBTOPIC CLUSTERS\\REPRESENTED}{Teal}\end{minipage}\hfill
\begin{minipage}[t]{0.190\linewidth}\metric{10}{INSTRUMENTATION /\\MODELLING RECORDS}{Sky}\end{minipage}\hfill
\begin{minipage}[t]{0.190\linewidth}\metric{8}{EFFECT ESTIMATES\\EXTRACTED}{Amber}\end{minipage}\hfill
\begin{minipage}[t]{0.190\linewidth}\metric{4}{REVIEWS /\\EVIDENCE SYNTHESES}{Clay}\end{minipage}

\vspace{6mm}

% -------------------------------------------------------------- signal of day
\section*{\sffamily\bfseries\color{Deep}Signal of the day}
\vspace{-1.5mm}{\color{Amber}\rule{\linewidth}{1.1pt}}\vspace{2.5mm}

\begin{keybox}
{\sffamily\bfseries\small\color{Deep}A US-derived PurpleAir correction does not travel.
Region-specific calibration is not a refinement --- it is a precondition.}\\[1.4mm]
\footnotesize
Ginsburg et al.\ (\textit{Aerosol Sci Technol}) evaluated three widely used US correction
equations against 1{,}013 PurpleAir units co-located within 500\,m of reference-grade
monitors on six continents (Jan 2023--Mar 2025). Raw readings overestimate PM$_{2.5}$ in
most regions and fail EPA sensor performance targets outright; the US-derived corrections
improve accuracy but \textbf{do not transfer outside North America}. A global quadratic
model helps in Australia/New Zealand and degrades elsewhere. Only regional calibrations
consistently met EPA 24-h thresholds, raising 24-h AQI category correctness to
$\geq$77\%. Read alongside Feng et al.\ (\textit{Meas Sci Technol}), whose review finds
no universally optimal strategy --- ML wins on predictive error
($<$10\,$\mu$g\,m$^{-3}$), physics-based models on interpretability, hybrids on
transferability --- and the operational conclusion for distributed networks is
uncomfortable: a network without local co-location has an unquantified bias, and no
published equation removes it. The mechanism is not mysterious. Both papers finger RH
sensitivity and aerosol-composition variability, which is exactly what a nephelometric
scattering-to-mass inversion cannot see.
\end{keybox}

\vspace{3mm}

% -------------------------------------------------------------- exec brief
\section*{\sffamily\bfseries\color{Deep}What changed today}
\vspace{-1.5mm}{\color{Amber}\rule{\linewidth}{1.1pt}}\vspace{2.5mm}

\footnotesize
\begin{itemize}
\item \textbf{Compound-hazard epidemiology arrives with adequate power.} Zhao et al.\
(\textit{JACC}) analyse 5.9 million CVD deaths across six countries and find co-occurring
heat waves and substantial fire-sourced air pollution carry \textbf{RR 1.115
(1.100--1.130)}, above heat-only (1.083) and smoke-only (1.014), with a small but
positive additive interaction (RERI 0.018, 95\% CI 0.002--0.035). The interaction term is
real but modest --- most of the compound risk is heat. Resist the synergy headline.

\item \textbf{Two independent imaging cohorts converge on subclinical coronary
structure.} Belay et al.\ tie cumulative bushfire-related PM$_{2.5}$ to total plaque
volume (\textbf{+10.14\,mm$^3$ per 1\,$\mu$g\,m$^{-3}$}, 95\% CI 1.45--18.82) in 365
serial CCTA participants at a 10-year mean of only $\sim$5.9\,$\mu$g\,m$^{-3}$; Chen
et al.\ tie 1-year PM$_{2.5}$/PM$_{10}$/BC to pericoronary fat attenuation index in
5{,}528 patients (+7.57/8.72/2.35\,HU per IQR in the RCA). Structural and inflammatory
endpoints, same direction, different populations.

\item \textbf{A well-designed PM$_{2.5}$ null.} In 1{,}587 ECHO-PATHWAYS children,
infancy PM$_{2.5}$ and NO$_2$ showed \textbf{no association} with atopic dermatitis at
4--6 years; only O$_3$ did (aRR 1.07 per 1\,ppb, 1.01--1.13), and only in children of
lower-educated mothers. Well-powered nulls from validated spatiotemporal models deserve
the same shelf space as positives.

\item \textbf{The satellite-fusion literature has started auditing its own bias.} Ding
et al.\ argue that ground-station clustering in cities systematically inflates retrieved
rural PM$_{2.5}$; their urban--rural balanced framework cuts clean-site RMSE from 18.18
to 8.31\,$\mu$g\,m$^{-3}$ and lowers the national mean by $\sim$55\%, which would
substantially revise PM-attributable mortality in low-density areas. If that holds, a
large body of exposure--response work inherits an upward exposure bias.

\item \textbf{Sensor siting became a tractable optimisation.} Patel et al.\
(\textit{AAAI}) maximise mutual information over continuous sensor coordinates by
gradient descent, with runtime independent of both placement count and candidate-grid
density --- it approaches greedy MI quality orders of magnitude faster. Directly usable
for network expansion planning.

\item \textbf{Caveat on a headline calibration number.} Msigwa et al.\ report RMSE
0.325\,$\mu$g\,m$^{-3}$ and $R^2$ 0.999 for a transformer/TCN hybrid against a reference
instrument. That is below the reference's own uncertainty and is not credible as a
generalisation estimate; it points to temporally adjacent train/test splits rather than
held-out deployment. Treat as an architecture description, not a performance claim.
\end{itemize}

\clearpage

% ==================================================================== ANALYTICS
\section*{\sffamily\bfseries\color{Deep}Corpus analytics}
\vspace{-1.5mm}{\color{Amber}\rule{\linewidth}{1.1pt}}\vspace{3mm}

\noindent\includegraphics[width=\linewidth]{\FIGDIR/f1_subtopics.png}
\figcap{Subtopic assignment of the 25 in-scope records. \textit{Sensing, forecasting \&
instrumentation} is today's largest cluster (6) --- an artefact of methodology, not of the
field: PubMed indexed exactly one instrumentation record in the 29 July entry window, and
five more arrive from the Consensus sweep that exists to cover journals PubMed does not
index. Read the instrumentation count as a backfill, not a daily flux.}

\vspace{2mm}
\noindent\includegraphics[width=\linewidth]{\FIGDIR/f2_design_endpoint.png}
\figcap{Left --- study architecture. Modelling/inventory and observational-cohort work tie
at 5 records each; reviews are 4/25 (16\%), lower than the 24\% of 27 July.
Right --- health endpoint. Nine records carry no human health endpoint at all --- five
sensing, three exposure-modelling and the risk-based AQI paper. That gap is the structural
feature of this issue: instrumentation advances and health inference are still being
published into separate literatures with no shared validation target.}

\vspace{2mm}
\noindent\includegraphics[width=\linewidth]{\FIGDIR/f3_metric_geography.png}
\figcap{Left --- particle metric. PM$_{2.5}$ alone is the exposure in 14 records, jointly
with PM$_{10}$ in 5; ultrafine particles appear once (an ecological country-level
analysis, the weakest design in the corpus), continuing the mismatch between deposition
physics and measurement practice. Right --- geographic provenance. China contributes 7;
the Global/multi-region block (10) is inflated by the instrumentation backfill, whose
papers are method-general. South Asia and sub-Saharan Africa contribute one record each.}

\vspace{2mm}
\noindent\includegraphics[width=\linewidth]{\FIGDIR/f6_lifecourse.png}
\figcap{Life-course windows addressed. Counts are non-exclusive. Adulthood/ageing
dominates (11) on the strength of the coronary-imaging and occupational records; the
instrumentation block (9) carries no specific window and is plotted for completeness.
Only two records address the in-utero window today, against five on 27 July.}

\vspace{2mm}
\noindent\includegraphics[width=\linewidth]{\FIGDIR/f4_heatmap.png}
\figcap{Subtopic $\times$ study-architecture cross-tabulation. The diagnostic cell is the
intersection of the \textit{Sensing} and \textit{Exposure assessment} rows with any
observational-epidemiology column: it is empty. Mechanistic toxicology is again wholly
experimental with no human biomarker bridge. Both patterns have now held across all three
issues in the corpus.}

\vspace{2mm}
\noindent\includegraphics[width=\linewidth]{\FIGDIR/f5_forest.png}
\figcap{All quantitative effect estimates in today's corpus on a common logarithmic axis.
\textbf{Exposure contrasts are not mutually comparable} and the plot shows direction and
precision only. The three Zhao et al.\ estimates share a common referent (no-hazard days)
and are comparable to each other --- that triplet is the informative comparison here. The
two ECHO-PATHWAYS estimates are per 1\,ppb O$_3$, retained because the accompanying
PM$_{2.5}$ null is itself the finding. Three Mokoena et al.\ odds ratios are plotted
without whiskers because the source reports no confidence intervals; they are
self-reported cross-sectional associations and are the weakest evidence shown.}

\clearpage

% ==================================================================== DIGEST
\section*{\sffamily\bfseries\color{Deep}Paper-level digest}
\vspace{-1.5mm}{\color{Amber}\rule{\linewidth}{1.1pt}}\vspace{1mm}

{\fontsize{7.4}{9}\selectfont\color{Slate}Ordered by cluster. Each entry gives the
methodological core and the finding that matters, not an abstract paraphrase.\par}

\band{Sensing, forecasting \& instrumentation\hfill 6 records}{Amber}

\paperentry{Ginsburg et al.\ 2026}
{Global and regional calibration of low-cost PM$_{2.5}$ sensors (PurpleAir)}
{Aerosol Sci Technol}
{1{,}013 PurpleAir units co-located within 500\,m of reference monitors on six
continents, Jan 2023--Mar 2025, assessed at 1-h and 24-h averaging by $R^2$, RMSE, nRMSE,
precision (CV) and AQI category correctness. Raw readings overestimate PM$_{2.5}$ in most
regions and miss EPA performance targets; US-derived corrections show limited
transferability outside North America; a data-driven global quadratic helps in
Australia/NZ and hurts elsewhere; \textbf{regional equations meet EPA 24-h thresholds in
multiple regions and lift 24-h AQI correctness to $\geq$77\%}. Limitation: co-location
proximity ($\leq$500\,m) is a proxy for comparability, so some regional ``bias'' is
spatial-gradient error, and continents are represented very unevenly. Still the strongest
available empirical case that LCS network calibration must be local.}
{10.1080/02786826.2026.2676293}

\paperentry{Feng et al.\ 2026}
{Calibration strategies and measurement frameworks for low-cost PM sensors}
{Meas Sci Technol}
{Review spanning ambient, mobile, indoor and workplace deployments (community networks,
vehicle-mounted platforms, smart-manufacturing floors, mining). Compares purely
data-driven, purely physics-based and hybrid physics-informed ML calibration on accuracy,
transferability, compute cost and interpretability: non-linear models reach
sub-10\,$\mu$g\,m$^{-3}$ error but resist embedded deployment and depend on
representative training data; \textbf{no strategy is universally optimal}. Names the
persistent barriers plainly --- sensor fouling, data governance, model generalisation,
hardware limits --- and points to self-cleaning optics, federated learning and
knowledge-distilled on-device calibration. Narrative rather than quantitative synthesis,
with no pooled performance estimate, so it frames rather than settles the question.}
{10.1088/1361-6501/ae6298}

\paperentry{Patton et al.\ 2026}
{Human health risk of particle exposure in outdoor environments, southwestern Pennsylvania}
{Environ Sci Pollut Res}
{ROCIS low-cost monitoring across 290 homes in Allegheny County, 2016--2021; optical
particle counts $>$0.5\,$\mu$m converted to PM$_{2.5}$-equivalent mass, census tracts
classified by environmental-justice status and industrial proximity, GEE models plus
non-cancer hazard quotients. \textbf{33.8\% of sites exceeded the WHO 24-h
15\,$\mu$g\,m$^{-3}$ guideline and 66.6\% the EPA annual 9\,$\mu$g\,m$^{-3}$ reference}
on 3-week site medians; winter concentrations significantly elevated. The count-to-mass
conversion is the weak link --- it assumes a fixed density and size distribution that
industrial plumes violate --- and the authors correctly flag the Sewickley extreme
(median 19.3\,$\mu$g\,m$^{-3}$) as resting on three homes. Directly relevant as a
template for community-scale LCS deployment with explicit EJ stratification.}
{10.1007/s11356-026-38083-2}

\paperentry{Patel et al.\ 2026}
{Scalable air-quality sensor placement via gradient-based mutual information maximization}
{Proc AAAI}
{Treats sensor coordinates as differentiable parameters and maximises mutual information
by continuous optimisation, so \textbf{runtime is independent of both the number of
placements and candidate-grid density} --- the two terms that make greedy MI intractable
at national scale. Beats random placement and the maximum-predictive-variance heuristic
on a continental PM$_{2.5}$ dataset and approaches greedy MI regionally, orders of
magnitude faster. Caveat: MI is computed under an assumed GP covariance, so placements
inherit that model's stationarity assumptions, which is exactly where real PM fields
misbehave near sources. Immediately useful for network-expansion planning under a fixed
sensor budget.}
{10.1609/aaai.v40i46.41257}

\paperentry{Huda et al.\ 2026}
{Low-cost sensor network for wildfire monitoring at remote community locations}
{J Air Waste Manag Assoc}
{Solar-powered microsystem (PM$_{2.5}$, CO, CO$_2$, O$_3$, tVOC) with a microcontroller
scheduler and mesh radio, built for forest deployment in northwestern Alberta where there
is no cellular or Wi-Fi coverage. Reported PM$_{2.5}$ RMSE $<$5\,$\mu$g\,m$^{-3}$; mesh
range validated at 1--5\,km through urban structures and vegetated terrain;
50$\times$37\,cm panel sized for treetop mounting with energy-aware duty cycling. This is
an engineering report, not an evaluation: no co-location period, drift characterisation
or seasonal uptime statistics are given, and winter operation is described as
prospective. Useful for the power and comms architecture rather than the data-quality
claim.}
{10.1080/10962247.2026.2665375}

\paperentry{Msigwa et al.\ 2026}
{AIoT real-time sensing and calibration of low-cost PM sensors (TCM-Net)}
{IEEE Internet Things J}
{Customised PMS7003 units with externally PWM-controlled fans to stabilise airflow, on a
oneM2M-compliant IoT platform, calibrated by a transformer + temporal-convolutional
network with a memory-aware module and trained against a co-located certified reference.
Reported RMSE 0.325\,$\mu$g\,m$^{-3}$, $R^2$ 0.999. \textbf{Those figures sit below the
reference instrument's own measurement uncertainty} and cannot be a valid generalisation
estimate; the near-certain explanation is temporally adjacent train/test splits on
autocorrelated series rather than held-out deployment. The forced-airflow hardware idea
is worth taking; the performance number is not.}
{10.1109/jiot.2025.3631631}

\band{Exposure assessment \& modelling\hfill 4 records}{Amber}

\paperentry{Ding et al.\ 2026}
{Urban--rural balanced estimation framework for PM$_{2.5}$ exposure across China}
{Geo-spat Inf Sci}
{Targets a structural bias rather than a model-fit problem: because ground stations
cluster in cities, AOD-to-PM$_{2.5}$ retrievals systematically overestimate rural
concentrations. Integrates AOD, reanalysis fields and ground measurements under
data-assimilation principles, adaptively constraining non-urban estimates with
low-weighted reanalysis samples, then ensemble-averages across dynamically weighted
background/observation settings. China 2015--2021: \textbf{clean-site RMSE 18.18
$\rightarrow$ 8.31\,$\mu$g\,m$^{-3}$, national mean PM$_{2.5}$ down $\sim$55\%}, with
corresponding correction of PM-attributable mortality in low-density areas. The 55\%
figure is the paper's own counterfactual, validated at clean sites that are still station
locations --- true rural accuracy remains unvalidated by construction. If even partly
right, it implies an inherited upward exposure bias across the satellite-based
exposure--response literature.}
{10.1080/10095020.2026.2624335}

\paperentry{Hu et al.\ 2026}
{Fusion of ground monitoring, satellite retrieval and CTM for seamless hourly 1-km PM$_{2.5}$}
{Geo-spat Inf Sci}
{Fills satellite AOD gaps with continuous WRF-Chem AOD, injects physics-informed features
(PM$_{2.5}$ temporal components, AOD vertical structure) from the CTM, then predicts with
a geospatially aware LightGBM that exploits monitoring-site spatial structure. Site-based
cross-validation $R^2$ 0.87, RMSE 13.42\,$\mu$g\,m$^{-3}$, MAE 8.68\,$\mu$g\,m$^{-3}$,
with day-and-night coverage demonstrated on one pollution episode. Site-based CV
overstates performance away from stations --- the standard weakness --- and an RMSE of
13.4 is large relative to the WHO annual guideline, so this is useful for episode mapping
and weak for long-term exposure assignment. The physics-feature injection is the
transferable step.}
{10.1080/10095020.2026.2680855}

\paperentry{Sanchez et al.\ 2026}
{Morbidity related to the level-two eruption of Popocatepetl}
{Indian J Occup Environ Med}
{Ecological correlational study in communities 15\,km from the crater, May 2023, with
daily PM$_{10}$/PM$_{2.5}$ and Kriging used to interpolate ash dispersion because the
monitoring network is too sparse to resolve it. 14\% of presentations were ash-related,
predominantly non-infectious respiratory; Spearman correlations between ash exposure and
morbidity were not statistically significant. The authors read this as possible
biological adaptation, which the design cannot support --- a null in a small ecological
study with interpolated exposure is far more likely to be measurement error and low
power. Its real value here is as a candid worked example of exposure assignment where no
formal network exists.}
{10.4103/ijoem.ijoem_146_24}

\paperentry{Nguyen et al.\ 2026}
{Satellite-driven PM$_{2.5}$ prediction: machine learning and explainable AI}
{Eng Res Express}
{Systematic review of 2015--2025 work integrating satellite products (LEO and
geostationary AOD, top-of-atmosphere reflectance, {\AA}ngstr\"om exponent) with ML/DL.
Finds Random Forest still a robust benchmark, hybrid and transformer architectures better
at spatiotemporal dependence, and AOD gap-filling plus normalisation the dominant
determinants of model stability; positions XAI as the route to physical plausibility. No
pooled quantitative comparison and no protocol registration, so the ranking claims are
narrative. Adds little beyond Hu and Ding above; retained as a citation anchor for the
physics-guided-learning argument.}
{10.1088/2631-8695/ae7028}

\clearpage

\band{Cardiovascular \& metabolic\hfill 4 records}{Teal}

\paperentry{Zhao et al.\ 2026}
{Compound extreme heat and fire-sourced air pollution vs.\ cardiovascular mortality}
{J Am Coll Cardiol}
{Multi-country time-series across six countries, 5.9 million CVD deaths, with each day
classified into mutually exclusive states (no hazard / heat-wave only / substantial
fire-sourced air pollution only / compound) and quasi-Poisson RRs plus formal additive
interaction. \textbf{Compound RR 1.115 (1.100--1.130)} vs.\ heat-only 1.083
(1.074--1.091) and SFAP-only 1.014 (1.008--1.021); RERI 0.018 (0.002--0.035),
attributable proportion 1.598\% (0.181--3.129). The interaction is statistically present
but small --- most compound risk is the heat term --- and the mutually-exclusive day
classification discards dose information within each state. Fire-sourced attribution
rests on modelled smoke fields, so exposure misclassification is non-differential but
real. Sets the design standard for compound-hazard work.}
{10.1016/j.jacc.2026.06.005}

\paperentry{Belay et al.\ 2026}
{Long-term bushfire-related PM$_{2.5}$ and coronary plaque burden}
{Eur J Prev Cardiol}
{365 asymptomatic adults aged 40--70 with a family history of CAD, serial coronary CTA at
baseline and 3 years, with daily total and bushfire-attributed PM$_{2.5}$ from validated
5\,km models assigned by residential postcode. Per 1\,$\mu$g\,m$^{-3}$ cumulative
bushfire PM$_{2.5}$: \textbf{total plaque volume +10.14\,mm$^3$ (1.45--18.82)} at
baseline and +4.92\,mm$^3$ (1.33--8.52) at follow-up; total PM$_{2.5}$ associations were
weaker. Note the exposure setting --- 10-year means of 5.88 and 5.93\,$\mu$g\,m$^{-3}$,
i.e.\ near the WHO guideline --- which makes a plaque-volume signal notable but also
means the bushfire and total metrics are highly collinear at these levels. Postcode-level
assignment in a 365-person sample is the binding limitation.}
{10.1093/eurjpc/zwag399}

\paperentry{Chen et al.\ 2026}
{Long-term PM exposure and coronary inflammation by pericoronary fat attenuation index}
{Environ Pollut}
{5{,}528 CCTA patients with a 656-person serial subcohort; FAI extracted for LAD, LCx and
RCA by a deep-learning pipeline; individual PM$_{2.5}$, PM$_{10}$ and black carbon from a
validated 1\,km spatiotemporal model; GAM and linear mixed-effects models. Per IQR of
1-year exposure, \textbf{RCA FAI rose 7.57 (PM$_{2.5}$), 8.72 (PM$_{10}$) and 2.35\,HU
(BC)}, with accelerated annual FAI progression longitudinally. FAI is technically
sensitive to tube voltage, contrast timing and epicardial adiposity, and a retrospective
clinical population undergoing CCTA is selected on suspected disease --- so absolute HU
shifts should not be read as population effects. The local-inflammation endpoint
nonetheless complements Belay's structural one.}
{10.1016/j.envpol.2026.128840}

\paperentry{Rusconi et al.\ 2026}
{Long-term ultrafine particle exposure 2010--2019 and cardiovascular burden in high-income countries}
{GeoHealth}
{Applies a new global UFP dataset to compute population-weighted exposure for high-income
countries and correlates it with country-level CVD DALYs, YLLs and YLDs, reporting
``consistent patterns'' between long-term UFP exposure and CVD burden. This is a
country-level ecological correlation with no individual data, no lag structure and no
control for the PM$_{2.5}$ with which UFP is strongly collinear --- the authors say so,
restrict to high-income countries deliberately, and warn against generalisation. It
cannot support causal or even population-attributable inference. Included only because it
is the corpus's sole ultrafine-focused record, which is itself the point: UFP
epidemiology still lacks the exposure infrastructure PM$_{2.5}$ has.}
{10.1029/2025GH001636}

\band{Respiratory \& allergic\hfill 3 records}{Sage}

\paperentry{Zhang et al.\ 2026}
{Infancy air pollution exposure and atopic dermatitis at 4--6 years}
{Epidemiology}
{Pooled longitudinal analysis of 1{,}587 children across six US cities in the
ECHO-PATHWAYS Consortium (2017--2023); first-year-of-life PM$_{2.5}$, NO$_2$ and O$_3$
from validated national spatiotemporal models; three caregiver-reported outcomes; Poisson
models with prespecified modification by child sex and maternal education.
\textbf{PM$_{2.5}$ and NO$_2$ null}; O$_3$ associated with current AD (aRR 1.07 per
1\,ppb, 1.01--1.13) and location-specific AD (1.06, 0.99--1.13), evident only in children
of lower-educated mothers. Outcome ascertainment is caregiver questionnaire, not clinical
examination, and the education-stratified finding is a subgroup result on a modest
sample. A well-powered PM$_{2.5}$ null against a validated exposure model is worth as
much as a positive.}
{10.1097/EDE.0000000000002033}

\paperentry{Li et al.\ 2026}
{Host genetic variation, airway microbiome and environmental exposure vs.\ respiratory health}
{mSystems}
{Sputum metagenomes from 1{,}651 individuals, with Mendelian randomisation implicating
SO$_2$, CO and PM$_{2.5}$ effects on lung function potentially mediated by
\textit{Actinomyces}, \textit{Haemophilus influenzae} and \textit{Veillonella} spp., and
genotype-dependent exposure--microbiome interactions (\textit{MEOX1} with lower
\textit{Filifactor alocis} under pollutant exposure; \textit{FAM110D} and \textit{USP36}
with higher \textit{Ralstonia pickettii}/\textit{Neisseria} and lower
\textit{Oribacterium}, tracking COPD risk). Fungal loci also reported. This is many
hypotheses from one cohort: MR on aggregate exposures with microbial mediators is fragile
to pleiotropy and the variant identifiers are rare or novel, so the loci need replication
before they mean anything. Mechanistically suggestive for why identical exposure yields
divergent lung-function trajectories.}
{10.1128/msystems.00442-26}

\paperentry{Mokoena et al.\ 2026}
{Socio-economic conditions, air pollution sources and eczema symptoms in Soshanguve teenagers}
{Public Health Chall}
{Cross-sectional modified ISAAC III survey of 2{,}885 schoolchildren aged 13--14 across
13 schools in Tshwane, South Africa; multiple logistic regression. Current eczema
symptoms 26.8\%, eczema ever 23.3\%. Frequent truck traffic near the home: AOR 1.26
(current) and 1.24 (ever); ETS 16--20 days/month: AOR 1.96; female sex, monthly
paracetamol use (AOR 2.37) and pet ownership also associated. No confidence intervals are
reported for most estimates, exposure is entirely self-reported with no measurement, and
the paracetamol association is the textbook reverse-causation trap. Value is geographic
coverage --- one of very few sub-Saharan records in this corpus --- not effect
estimation.}
{10.1002/puh2.70325}

\clearpage

\band{Neuro, cognition \& mental health\hfill 2 records}{Deep}

\paperentry{Kabir et al.\ 2026}
{Environmental attributes as moderators of neighbourhood deprivation and mental health}
{Health Place}
{Systematic review across seven databases plus grey literature, 19 studies included,
asking specifically whether objectively measured environmental attributes moderate the
deprivation--mental-health association. Moderation findings were mixed: greenness and
blue space generally protective in disadvantaged neighbourhoods, air pollution
exacerbated poor mental health, parks and walkability inconsistent. Noise and
elder-support centres were each examined in exactly one study. With 19 heterogeneous
studies and no meta-analysis, the moderation claim is directional at best, and the
authors' call for longitudinal designs is the honest conclusion. Relevant if you intend
to use LCS network output as an effect modifier rather than a main exposure.}
{10.1016/j.healthplace.2026.103719}

\paperentry{Sarkar et al.\ 2026}
{PM-induced neurotoxicity: neuroinflammation and organelle dysfunction}
{J Appl Toxicol}
{Narrative review mapping CNS entry routes (olfactory-bulb translocation, blood--brain
barrier disruption, peripheral inflammatory signalling) onto subcellular failure modes:
unfolded protein response with CHOP and caspase-12 induction, impaired oxidative
phosphorylation, ROS accumulation, mitochondrial permeability transition, and compromised
autophagic flux and lysosomal integrity. Coherent as a mechanistic map but adds no new
data, and --- as with most of this literature --- draws heavily on instillation and
high-concentration in vitro models whose dosimetry does not correspond to ambient
exposure. Use for the organelle-level framing, not for dose--response.}
{10.1002/jat.70360}

\band{Mechanistic toxicology\hfill 2 records}{Violet}

\paperentry{Wang et al.\ 2026}
{PM exposure, mitochondrial stress and ECM remodelling in scleral fibroblasts}
{Exp Eye Res}
{Combined in vivo and in vitro: topically applied PM produced particle-surrogate
distribution to ocular tissues with scleral thinning and axial elongation, while human
scleral fibroblasts showed raised oxidative stress, mitochondrial ultrastructural
abnormality, phenotypic change and altered collagen metabolism; transcriptomics
implicated ECM organisation, mitochondrial metabolism and stress-response pathways. The
framing as a myopia mechanism is plausible but the exposure route is the problem ---
topical application is not inhalation, and axial elongation in a rodent eye is a weak
proxy for human myopia progression. Notable mainly as one of very few records addressing
an ocular endpoint at all.}
{10.1016/j.exer.2026.111187}

\paperentry{Liu et al.\ 2026}
{Resveratrol ameliorates PM$_{2.5}$-aggravated IBS-D via microbiota and metabolite restoration}
{Environ Toxicol Pharmacol}
{Three-arm murine design (IBS-D, PM$_{2.5}$+IBS-D, resveratrol+PM$_{2.5}$+IBS-D) with
untargeted metabolomics, 16S microbiota profiling, colonic histopathology,
oxidative-stress and barrier-function assays and visceral-hypersensitivity testing.
PM$_{2.5}$ worsened IBS-D with dysbiosis and disrupted hypoxanthine and linoleic-acid
metabolism; resveratrol restored both. No exposure concentrations, particle source or
dose--response are reported in the abstract, group sizes are unstated, and the
resveratrol arm is a single dose --- so this is proof-of-concept, not a therapeutic
result. The gut-axis endpoint remains underexplored, which is why it is retained.}
{10.1016/j.etap.2026.105120}

\band{Occupational \& indoor exposure\hfill 2 records}{Clay}

\paperentry{Qi et al.\ 2026}
{Emission rates of respirable crystalline silica and dust from engineered stone fabrication}
{Ann Work Expo Health}
{Extracts emission rates (mass per time) from four published experimental simulation
studies of cutting and grinding engineered stone, tests control technology and task as
determinants by regression, then drives a two-zone mathematical model and validates
predictions against task-duration field measurements. RCS emission rate rose with slab
crystalline-silica content and fell with control technology; \textbf{near-field
predictions tended to underestimate measured exposures} though the distributions
overlapped. Four source studies is a thin evidence base and the underestimation direction
is exactly the wrong way for a protective standard. Methodologically the most
transferable record in this cluster --- an explicit emission-rate-to-exposure model
rather than another exposure survey.}
{10.1093/annweh/wxag060}

\paperentry{Ghosh et al.\ 2026}
{Respiratory quality of life and pulmonary function in traffic police, eastern India}
{Indian J Occup Environ Med}
{Cross-sectional comparison of 62 traffic police against 120 unexposed controls
(Oct 2023--Jan 2024) using the St George's Respiratory Questionnaire and digital
spirometry. Exposed personnel had significantly higher SGRQ scores and lower FVC and
FEV$_1$ (both $p<0.001$), with SGRQ correlating positively with pulmonary function
measures. No exposure measurement whatsoever --- occupation is the exposure proxy --- and
no adjustment for smoking, duration of service or socioeconomic position is reported,
while simple random sampling from unequal groups invites selection differences. Reads as
a hazard-surveillance signal for a large occupational group, not as an exposure--response
result.}
{10.4103/ijoem.ijoem_150_25}

\band{Reproductive \& developmental\hfill 1 record}{Sky}

\paperentry{Lin et al.\ 2026}
{Maternal PM exposure during pregnancy and offspring liver function}
{Environ Pollut}
{Fujian Birth Cohort, 1{,}046 mother--offspring pairs, with gestational and
trimester-specific PM$_1$, PM$_{2.5}$ and PM$_{10}$, serum ALT and AST at ages 1--2
years, restricted cubic splines for shape and distributed lag models for timing.
\textbf{Second-trimester IQR increments gave ALT +1.42 (0.82--2.03), +1.41 (0.77--2.05)
and +0.91 (0.35--1.46)\,U/L} for PM$_1$/PM$_{2.5}$/PM$_{10}$, with AST increases of
similar magnitude. The effect sizes are small relative to the assay's biological
variability in toddlers and clinical significance is unestablished; a single post-natal
measurement precludes trajectory inference, and PM$_1$/PM$_{2.5}$/PM$_{10}$ are too
collinear to separate as cleanly as the paper implies. The distributed-lag identification
of a second-trimester window is the durable part.}
{10.1016/j.envpol.2026.128835}

\band{Burden, policy \& mitigation\hfill 1 record}{Coral}

\paperentry{Kraus et al.\ 2026}
{A risk-based air quality index for Germany}
{Int J Hyg Environ Health}
{Rebuilds the German AQI on risk equivalence rather than convention: pollutant--outcome
pairs with causal or likely-causal evidence were identified, effect estimates extracted
from systematic reviews, daily-mean effects transformed to hourly using German monitoring
data, and equivalence coefficients derived by Monte Carlo simulation standardised to
PM$_{2.5}$ as reference. PM$_{2.5}$ thresholds anchor to WHO 2021 and the EU information
threshold; PM$_{10}$, NO$_2$, O$_3$ and SO$_2$ thresholds follow by risk equivalence,
with an added multi-pollutant module. Evaluated on 2019 and 2022 hourly station data: a
marked shift from ``very good'' to ``good''/``moderate'' and modestly more ``poor'' days.
The daily-to-hourly effect transformation is the load-bearing and least verifiable
assumption, and the index inherits every bias of the source meta-analyses. The most
directly policy-relevant record today, and a clean template for making an index
defensible rather than traditional.}
{10.1016/j.ijheh.2026.114867}

\clearpage

% ==================================================================== TABLE
\section*{\sffamily\bfseries\color{Deep}Corpus register}
\vspace{-1.5mm}{\color{Amber}\rule{\linewidth}{1.1pt}}\vspace{2.5mm}

{\fontsize{6.7}{8.4}\selectfont
\setlength{\tabcolsep}{2.4pt}
\begin{longtable}{@{}>{\RaggedRight}p{25mm}>{\RaggedRight}p{41mm}>{\RaggedRight}p{27mm}%
>{\RaggedRight}p{23mm}>{\RaggedRight}p{18mm}>{\RaggedRight}p{15mm}@{}}
\toprule
\textsf{\textbf{First author}} & \textsf{\textbf{Focus}} & \textsf{\textbf{Journal}} &
\textsf{\textbf{Design}} & \textsf{\textbf{Metric}} & \textsf{\textbf{Region}}\\
\midrule
\endfirsthead
\toprule
\textsf{\textbf{First author}} & \textsf{\textbf{Focus}} & \textsf{\textbf{Journal}} &
\textsf{\textbf{Design}} & \textsf{\textbf{Metric}} & \textsf{\textbf{Region}}\\
\midrule
\endhead
\multicolumn{6}{@{}l@{}}{\cellcolor{Mist}\textsf{\textbf{\textcolor{Deep}{Neuro / mental health}}}}\\[0.6mm]
\href{https://doi.org/10.1001/jamanetworkopen.2026.25015}{Dales} & Wildfire-sourced PM$_{2.5}$ and ED visits for migraine and other primary headache & \textit{JAMA Netw Open} & Case-crossover & PM$_{2.5}$ (wildfire) & Canada\\
\href{https://doi.org/10.1093/ageing/afag215}{Glover} & Umbrella meta-review of urban environmental exposures and cognitive health & \textit{Age Ageing} & Umbrella review / meta-analysis & PM$_{2.5}$ & Global\\
\href{https://doi.org/10.1016/j.nsa.2026.107019}{Goossens} & In utero air pollution exposure and pre-/postnatal brain development & \textit{Neurosci Appl} & Narrative review & PM$_{2.5}$ / PM$_{10}$ / PM$_{\mathrm{coarse}}$ & Global\\
\href{https://doi.org/10.1016/j.isci.2026.116617}{Ma} & PM$_{2.5}$-induced neurotoxicity and depression: exposure risks and mechanisms & \textit{iScience} & Narrative review & PM$_{2.5}$ + components & Global\\
\href{https://doi.org/10.3760/cma.j.cn112150-20250804-00752}{Pan} & Maternal PM exposure and offspring neurodevelopment: effects and mechanisms & \textit{Zhonghua Yu Fang Yi Xue Za Zhi} & Narrative review & PM (bulk) & China\\
\addlinespace[0.9mm]
\multicolumn{6}{@{}l@{}}{\cellcolor{Mist}\textsf{\textbf{\textcolor{Teal}{Cardiovascular \& metabolic}}}}\\[0.6mm]
\href{https://doi.org/10.3760/cma.j.cn112338-20251120-00843}{Huang} & PM$_{2.5}$ and dyslipidaemia in rural Chinese children; dietary effect modification & \textit{Zhonghua Liu Xing Bing Xue Za Zhi} & Cross-sectional (multistage) & PM$_{2.5}$ & China\\
\href{https://doi.org/10.32394/pe/218966}{Kazmierska} & Air pollution and cardiovascular disease: review of current data & \textit{Przegl Epidemiol} & Narrative review & PM$_{2.5}$ / PM$_{10}$ & Global\\
\href{https://doi.org/10.1007/s00125-026-06814-2}{Xi} & Air pollution, built environment and temperature vs. prediabetes phenotypes & \textit{Diabetologia} & Prospective cohort & PM$_{2.5}$ & Germany\\
\href{https://doi.org/10.1161/JAHA.125.046917}{Yang} & Short-term PM$_{2.5}$ and hospitalisation/ED use in atrial fibrillation & \textit{J Am Heart Assoc} & Case-crossover & PM$_{2.5}$ (wildfire-influenced) & USA\\
\addlinespace[0.9mm]
\multicolumn{6}{@{}l@{}}{\cellcolor{Mist}\textsf{\textbf{\textcolor{Sky}{Reproductive \& developmental}}}}\\[0.6mm]
\href{https://doi.org/10.1016/j.jes.2025.10.043}{Chen} & Prenatal air pollution, fetal growth trajectory phenotypes and child neurodevelopment & \textit{J Environ Sci (China)} & Prospective cohort & PM$_{2.5}$ / PM$_{10}$ & China\\
\href{https://doi.org/10.1016/j.ecoenv.2026.120566}{Huang} & Joint prenatal bisphenol + air-pollutant mixture exposure and preschool IQ & \textit{Ecotoxicol Environ Saf} & Prospective cohort (mixtures) & PM$_{2.5}$ / PM$_{10}$ & China\\
\href{https://doi.org/10.1007/s40572-026-00552-8}{Peterson} & Ambient air pollution, PAHs and pubertal development: critical review & \textit{Curr Environ Health Rep} & Critical review & PM$_{2.5}$ components & Global\\
\href{https://doi.org/10.1016/j.envres.2026.125313}{Xie} & Long-term PM$_{2.5}$ and timing of pubertal development in boys and girls & \textit{Environ Res} & Prospective cohort & PM$_{2.5}$ & China\\
\addlinespace[0.9mm]
\multicolumn{6}{@{}l@{}}{\cellcolor{Mist}\textsf{\textbf{\textcolor{Sage}{Respiratory \& allergic}}}}\\[0.6mm]
\href{https://doi.org/10.1007/s00484-026-03289-z}{Bayar Muluk} & Climate change and otolaryngology: airway disease impacts and emission reduction & \textit{Int J Biometeorol} & Narrative review & PM$_{2.5}$ & Global\\
\href{https://doi.org/10.4103/mgr.MEDGASRES-D-25-00164}{Liao} & Short-term air pollution, respiratory symptoms and lung function in asthma & \textit{Med Gas Res} & Panel study & PM$_{2.5}$ / PM$_{10}$ & China\\
\href{https://doi.org/10.3390/epidemiologia7040097}{Sanchez de Prada} & Ambient PM and asthma-coded admissions in children 0-3 y, Spain & \textit{Epidemiologia} & Ecological panel & PM$_{2.5}$ / PM$_{10}$ & Spain\\
\href{https://doi.org/10.1016/j.jhazmat.2026.143076}{Yang} & Allergenic vs. toxic risk of airborne seed hairs compared with ambient PM$_{2.5}$ & \textit{J Hazard Mater} & Field measurement & PM$_{2.5}$ + bioaerosol & China\\
\addlinespace[0.9mm]
\multicolumn{6}{@{}l@{}}{\cellcolor{Mist}\textsf{\textbf{\textcolor{Violet}{Mechanistic toxicology}}}}\\[0.6mm]
\href{https://doi.org/10.1016/j.jes.2025.12.074}{Ding} & PM$_{2.5}$ promotes vascular calcification via miR-27b-5p/HES1 signalling & \textit{J Environ Sci (China)} & In vitro + organoid + murine & PM$_{2.5}$ & China\\
\href{https://doi.org/10.1007/s00360-026-01698-5}{Maraschi} & Osmoregulatory disruption and metal bioaccumulation from settleable atmospheric PM & \textit{J Comp Physiol B} & Controlled exposure (ecotox) & Settleable APM & Brazil\\
\href{https://doi.org/10.1016/j.jes.2025.11.018}{Wu} & Macrophage NLRP3 inflammasome in real-ambient PM-induced hepatic glucose dysregulation & \textit{J Environ Sci (China)} & Real-ambient murine + co-culture & Real-ambient PM & China\\
\addlinespace[0.9mm]
\multicolumn{6}{@{}l@{}}{\cellcolor{Mist}\textsf{\textbf{\textcolor{Clay}{Occupational \& indoor}}}}\\[0.6mm]
\href{https://doi.org/10.3389/fmedt.2026.1796828}{Chen} & Nested suction instrument to cut intraoperative surgical-smoke PM$_{2.5}$ & \textit{Front Med Technol} & Bench / device experiment & PM$_{2.5}$ (surgical smoke) & China\\
\href{https://doi.org/10.1093/annweh/wxag055}{Cole} & Determinants of respirable crystalline silica exposure in Australian tunnelling & \textit{Ann Work Expo Health} & Observational exposure study & Respirable dust / RCS & Australia\\
\href{https://doi.org/10.1016/j.jhazmat.2026.143081}{Ding} & Stage-resolved hair Zn-isotope dynamics as a firefighting combustion-exposure tracer & \textit{J Hazard Mater} & Repeated-measures biomarker & Combustion-derived PM & China\\
\href{https://doi.org/10.3390/jox16040119}{Ramsden} & Xenobiotic hazards in aircraft cabin air & \textit{J Xenobiot} & Review + risk economics & Ultrafine particles & Global\\
\addlinespace[0.9mm]
\multicolumn{6}{@{}l@{}}{\cellcolor{Mist}\textsf{\textbf{\textcolor{Amber}{Sensing, forecasting \& instrumentation}}}}\\[0.6mm]
\href{https://doi.org/10.1016/j.jenvman.2026.130515}{Shi} & Adaptive spatiotemporal graph transformer for multi-site multi-step PM$_{2.5}$ forecasting & \textit{J Environ Manage} & Deep-learning model & PM$_{2.5}$ & China / India\\
\href{https://doi.org/10.3390/bios16070380}{Tsuchiya} & Seasonal and indoor-outdoor characteristics of fluorescent aerosol particles in offices & \textit{Biosensors} & Real-time sensor campaign & Bioaerosol / coarse & Japan\\
\addlinespace[0.9mm]
\multicolumn{6}{@{}l@{}}{\cellcolor{Mist}\textsf{\textbf{\textcolor{Amber}{Exposure assessment \& modelling}}}}\\[0.6mm]
\href{https://doi.org/10.1029/2025GH001621}{Solorzano-Araque} & Multiscale exposure modelling: LOTOS-EUROS + mobility surveys + particle toxicity & \textit{GeoHealth} & CTM + mobility integration & PM$_{2.5}$ / PM$_{10}$ & Colombia\\
\addlinespace[0.9mm]
\multicolumn{6}{@{}l@{}}{\cellcolor{Mist}\textsf{\textbf{\textcolor{Coral}{Burden, policy \& mitigation}}}}\\[0.6mm]
\href{https://doi.org/10.1002/hsr2.72886}{Karamzad} & Ischaemic heart disease burden in MENA 1990-2021; PM attributable fraction 28.2\% & \textit{Health Sci Rep} & GBD secondary analysis & Ambient PM & MENA\\
\href{https://doi.org/10.1371/journal.pone.0347989}{Rees} & HEPA filtration to reduce respiratory infection in care homes & \textit{PLoS One} & RCT process evaluation & Filtered PM & UK\\
\href{https://doi.org/10.1016/j.jes.2026.01.009}{Shan} & Policy-driven emission reductions and post-pandemic rebound, high-resolution inventories & \textit{J Environ Sci (China)} & Emission inventory & PM emissions & China\\
\addlinespace[0.9mm]
\multicolumn{6}{@{}l@{}}{\cellcolor{Mist}\textsf{\textbf{\textcolor{Slate}{Other clinical endpoints}}}}\\[0.6mm]
\href{https://doi.org/10.3390/dj14070440}{Lupita} & PM$_{2.5}$, preventive-care capacity and caries/edentulism burden across the EU & \textit{Dent J} & Ecological panel & PM$_{2.5}$ & EU-27\\
\href{https://doi.org/10.1016/j.jdin.2026.06.008}{Park} & Long-term PM exposure and risk of androgenetic alopecia & \textit{JAAD Int} & Nationwide cohort & PM$_{2.5}$ / PM$_{10}$ & South Korea\\
\href{https://doi.org/10.3171/2026.3.PEDS2615}{Sun} & Environmental exposure and area deprivation vs. paediatric hydrocephalus outcomes & \textit{J Neurosurg Pediatr} & Registry cohort + geospatial & PM$_{2.5}$ (satellite-derived) & USA\\
\addlinespace[0.9mm]
\bottomrule
\end{longtable}}


\vspace{4mm}
\begin{keybox}
{\sffamily\bfseries\footnotesize\color{Deep}Provenance and method}\\[1.2mm]
{\fontsize{7.2}{8.8}\selectfont
\textbf{PubMed} (NCBI E-utilities) was queried on 29 July 2026 against
\texttt{[Date - Entry]} for 2026/07/29 on two separate axes: particle terms crossed with
health-outcome terms (18 hits) and monitoring, sensing and calibration terms crossed with
particle terms (1 hit, off-topic). \textbf{Europe PMC} \texttt{CREATION\_DATE} for the
same day returned 1 record, already present in the PubMed set. \textbf{OpenAlex} returned
HTTP 429/403 on the created-date filter --- that leg remains broken, as recorded on
28 July. The \textbf{arXiv} relevance query returned 60 entries, none submitted on or
after 20 July 2026, so it contributed nothing; its date-range parameter is still
unusable. Because PubMed structurally does not index \textit{Atmos.\ Meas.\ Tech.},
\textit{Aerosol Sci.\ Technol.}, \textit{Meas.\ Sci.\ Technol.}\ or the IoT and ML venues
where sensor-calibration work appears, eight instrumentation and modelling records were
added from a \textbf{Consensus} sweep (LCS calibration and drift, co-location against
reference monitors, satellite--surface fusion, network siting, edge inference). Consensus
offers no entry-date filter, so those eight are dated by DOI registration (2025-11 to
2026-06) rather than by a 29 July entry window, and are flagged as a backfill wherever
counts are reported --- notably in the subtopic figure. All eight were checked against
\texttt{state/seen.json} and none had appeared in a previous issue.
\textbf{3 records excluded}: a lithium-ion-battery ultrasonic-sensing paper
(PMID 42525745, matched on ``sensor'' + ``machine learning'', no PM content); a workshop
proceedings with no data (PMID 42524492); and one IoT paper whose abstract was
unavailable from every source (\texttt{10.1109/jiot.2026.3707055}) --- reasons logged to
\texttt{state/rejected.jsonl}.
\\[1.0mm]
\textbf{Correction, reissued 31 July 2026.} As first published, three records in this
issue carried DOIs that PubMed's source metadata had registered incorrectly, so those
links misresolved: Patton et al.\ (PMID 42521969) pointed to a 2019 \textit{Sci Rep}
article, Rusconi et al.\ (PMID 42524698) to an \textit{Environ Sci Technol} article, and
Mokoena et al.\ (PMID 42524488) to an \textit{IJERPH} article. All three have since been
corrected upstream and are now given as
\texttt{10.1007/s11356-026-38083-2}, \texttt{10.1029/2025GH001636} and
\texttt{10.1002/puh2.70325} respectively. No other content in this issue has changed.
\\[1.0mm]
\textbf{Clinical-trial watch}: one registration was created or updated in the
2026-07-23$\rightarrow$29 window --- NCT05874479, a 440-participant sham-controlled
portable air-cleaner trial for blood pressure in New York City public housing
(NYU Langone, recruiting, primary completion 2027-09). Accumulated to
\texttt{state/trials.json} and surfaced in the weekly rollup, not here.
Subtopic, design, particle-metric and geography labels are assigned from title, abstract,
keywords and MeSH terms; assignment is single-label by dominant contribution and is a
judgement, not a controlled vocabulary. Effect estimates are transcribed verbatim from
abstracts and use heterogeneous exposure contrasts --- see the forest-plot caption.
Content derived from publisher metadata; abstracts remain the copyright of their
respective publishers.\par}
\end{keybox}

\end{document}
